\documentclass[border=5pt]{standalone}
\usepackage{tikz}
\usetikzlibrary{calc}

\begin{document}
\begin{tikzpicture}
    % Define coordinates
    \coordinate (A) at (0, 5);
    \coordinate (D) at (0, 0);
    \coordinate (B) at (-4, 0);
    \coordinate (C) at (3, 0);
    
    % Rectangle height
    \def\recty{2}
    
    % Calculate E and H
    % Line AB: y = 1.25(x+4)
    % At y=2, x = -2.4
    \coordinate (E) at (-2.4, \recty);
    \coordinate (F) at (-2.4, 0);
    
    % Line AC: y = -5/3(x-3)
    % At y=2, x = 1.8
    \coordinate (H) at (1.8, \recty);
    \coordinate (G) at (1.8, 0);
    
    % M is intersection
    \coordinate (M) at (0, \recty);
    
    % Draw Triangle Sides
    \draw (A) -- (B);
    \draw (A) -- (C);
    
    % Draw Base Line segments
    % FB solid
    \draw (B) -- (F);
    % DF dashed (Instruction: change DF to dashed)
    \draw[dashed] (F) -- (D);
    % GD solid (Instruction: keep GD solid)
    \draw (D) -- (G);
    % GC solid
    \draw (G) -- (C);
    
    % Draw Altitude
    \draw (A) -- (D);
    
    % Draw Rectangle
    \draw (E) -- (F);
    \draw (H) -- (G);
    \draw (E) -- (H);
    
    % Right angle at M
    \draw ($(M) + (0.15, 0)$) -- ($(M) + (0.15, 0.15)$) -- ($(M) + (0, 0.15)$);
    
    % Labels
    \node[above] at (A) {A};
    \node[below left] at (B) {B};
    \node[below right] at (C) {C};
    \node[below] at (D) {D};
    \node[left] at (E) {E};
    \node[below] at (F) {F};
    \node[below] at (G) {G};
    \node[right] at (H) {H};
    \node[left] at (M) {M};
\end{tikzpicture}
\end{document}
